% JEPASafe mini-project report
% Compile:  pdflatex jepasafe_report.tex   (run twice for the ToC/refs)
%      or:  latexmk -pdf jepasafe_report.tex
%
% TODO before submission: fill in the author names and roll numbers below,
% and verify every reference in Section 10 against the real publication.

\documentclass[11pt,a4paper]{article}

\usepackage[utf8]{inputenc}
\usepackage[T1]{fontenc}
\usepackage[margin=1in]{geometry}
\usepackage{amsmath,amssymb}
\usepackage{graphicx}
\usepackage{booktabs}
\usepackage{array}
\usepackage{fancyhdr}
\usepackage{listings}
\usepackage{xcolor}
\usepackage[hidelinks]{hyperref}

\pagestyle{fancy}
\fancyhf{}
\lhead{Mini Project Report}
\rhead{FPGA Accelerator for Human Activity Recognition}
\cfoot{\thepage}
\renewcommand{\headrulewidth}{0.4pt}

\lstset{
  basicstyle=\ttfamily\footnotesize,
  breaklines=true,
  frame=single,
  framesep=4pt,
  rulecolor=\color{gray!50},
  columns=fullflexible,
}

\title{\textbf{Design and FPGA Implementation of a 1D~CNN--LSTM\\
Accelerator for Real-Time Human Activity Recognition}}

% TODO: replace with the actual team
\author{Anubhav (ROLL\_NO) \\
National Institute of Technology, Tiruchirappalli}

\date{October 18, 2026}

\begin{document}

\maketitle

\begin{abstract}
This project presents the design, hardware translation and resource
characterisation of a compact 1D~convolutional--recurrent neural network for
human activity recognition from inertial sensor data, targeting a Xilinx
Zynq~UltraScale+ MPSoC. A 1D~CNN front end followed by an LSTM classifier is
trained on the UCI~HAR dataset and reaches 90.91\% test accuracy using only
4{,}374 parameters across six sensor channels. The trained network is converted
to synthesisable C++ using \texttt{hls4ml}, compiled to register-transfer level
by Vitis~HLS and packaged as an AXI4-Stream IP core. Post-synthesis
characterisation on the XCZU7EV device reports 33{,}500~LUT (14\% of the part),
54~DSP slices (3\%), 159 block~RAMs (25\%) and an inference latency of
1.60\,ms at 100\,MHz, against a 2.56\,s input window --- a margin of three
orders of magnitude. The work additionally quantifies how high-level synthesis
configuration, rather than network size, dominates hardware cost: a single
change of activation layer construct reduced the two activation blocks from
222{,}550~LUT to 2{,}182~LUT, and selecting a streaming interface over a
parallel one reduced the complete design from approximately 298{,}000~LUT to
35{,}745~LUT. The quantization strategy is documented in full, including a derivation of
\texttt{hls4ml}'s automatic accumulator widening that reproduces all four
generated accumulator types exactly. A negative result is also reported: the
design does not fit a low-cost Gowin GW1NR-9 device, because the LSTM's sigmoid and tanh lookup
tables alone require more block~RAM than that device contains. Results in this
report are high-level-synthesis estimates; bitstream generation and on-board
deployment remain future work.
\end{abstract}

\tableofcontents
\newpage

\section{Introduction}

Human activity recognition (HAR) from body-worn inertial sensors underpins
applications in healthcare monitoring, fall detection and personal safety
devices. Such systems are typically battery-powered and latency-sensitive,
which motivates performing inference at the edge rather than transmitting raw
sensor streams to a server.

Field-programmable gate arrays (FPGAs) are attractive for this role. Unlike a
general-purpose processor, an FPGA executes a neural network as dedicated
parallel hardware, and unlike a fixed-function ASIC it can be reprogrammed as
the model evolves. The difficulty is that translating a trained network into
efficient hardware is not automatic: the same network can differ in area by two
orders of magnitude depending on how it is mapped.

This project builds the complete translation chain --- train, convert,
synthesise, package --- for a 1D~CNN--LSTM activity classifier, and measures
what each stage actually costs in hardware. The specific contributions are:

\begin{enumerate}
  \item A compact CNN--LSTM achieving 90.91\% six-class accuracy on UCI~HAR
        using 4{,}374 parameters and only the six sensor channels an MPU6050
        physically provides.
  \item A complete, reproducible \texttt{hls4ml}~$\rightarrow$~Vitis~HLS
        $\rightarrow$ Vivado~IP flow, with the generated core's interface
        characterised at the port level.
  \item Quantitative resource characterisation across interface type,
        arithmetic precision, activation-table size and network width,
        isolating which choices matter and which do not.
  \item A documented negative result establishing why gated recurrence is
        infeasible on a small ($<$10k~LUT) FPGA, with the block~RAM accounting
        that explains it.
\end{enumerate}

\section{Related Work}

\subsection{Deep learning for sensor-based activity recognition}

The combination of convolutional feature extraction with recurrent temporal
modelling is the established architecture for wearable HAR. Ordóñez and
Roggen~\cite{ordonez2016} introduced DeepConvLSTM, stacking 1D~convolutional
layers beneath LSTM layers, and demonstrated that the convolutional front end
learns local motion primitives while the recurrent stage captures their
temporal ordering. Hammerla et al.~\cite{hammerla2016} systematically compared
deep, convolutional and recurrent models across wearable HAR benchmarks and
found recurrent architectures particularly effective for the naturalistic,
variable-duration activities typical of this domain.

The architecture used here follows that lineage directly, but at drastically
reduced width. Where DeepConvLSTM employs 64 convolutional filters and
128~hidden units, this work targets 16 of each, because the parameter budget is
set by on-chip memory rather than by accuracy alone.

\subsection{The UCI HAR benchmark}

Anguita et al.~\cite{anguita2013} published the UCI~HAR dataset, comprising
recordings from 30~participants performing six activities with a waist-mounted
smartphone sampled at 50\,Hz. Its wide adoption provides published baselines for
comparison, which is the reason it was selected here over sensor-matched
alternatives. Its principal limitation for this project is discussed in
Section~\ref{sec:limitations}: it contains no fall or distress class.

\subsection{Neural network inference on FPGAs}

Two toolflows dominate academic FPGA inference. FINN, introduced by Umuroglu et
al.~\cite{umuroglu2017}, targets heavily quantised (binarised) networks and
generates streaming dataflow architectures. \texttt{hls4ml}, introduced by
Duarte et al.~\cite{duarte2018} for microsecond-latency inference in particle
physics triggers, converts Keras and PyTorch models into Vivado/Vitis~HLS C++
and exposes per-layer control over fixed-point precision and multiplier reuse.

This project uses \texttt{hls4ml} because its reuse-factor abstraction maps
cleanly onto the design question here: the available latency budget (2.56\,s) is
enormous relative to the computation, so the appropriate strategy is to trade
all of that latency for reduced area. Section~\ref{sec:results} shows this
trade is worth roughly an order of magnitude.

\subsection{Joint-embedding predictive architectures}

The project was initially scoped around a Transformer-based joint-embedding
predictive architecture (JEPA), following LeCun's
proposal~\cite{lecun2022} and the image-domain realisation of Assran et
al.~\cite{assran2023}. That direction was abandoned early: self-attention and
layer normalisation are expensive to synthesise, and the self-supervised
pre-training that gives JEPA its advantage requires far more unlabelled data
than this project could collect. The supervised CNN--LSTM was adopted instead.
The project name retains the original acronym for continuity with earlier
documentation.

\section{Theoretical Background}

\subsection{Inertial sensing and windowing}

A six-axis inertial measurement unit reports linear acceleration along three
orthogonal axes and angular velocity about those same axes. Activity is not
identifiable from an instantaneous sample; it is a property of a time interval.
The signal is therefore segmented into fixed-length windows. At the UCI~HAR
sampling rate of 50\,Hz, the 128-sample window used here spans

\begin{equation}
T_{\text{window}} = \frac{128\ \text{samples}}{50\ \text{Hz}} = 2.56\ \text{s},
\end{equation}

long enough to contain several gait cycles. This figure sets the real-time
constraint: inference must complete within 2.56\,s to keep pace with the sensor.

\subsection{1D convolution for time series}

A 1D convolutional layer slides $F$ learned kernels of length $k$ along the time
axis, each spanning all $C$ input channels. With no padding, an input of $T$
time steps produces

\begin{equation}
T_{\text{out}} = T - k + 1
\end{equation}

output steps and requires

\begin{equation}
\text{MACs} = T_{\text{out}} \times k \times C \times F
\end{equation}

multiply--accumulate operations. This expression matters for hardware: the MAC
count, not the parameter count, determines arithmetic cost, and the factor
$T_{\text{out}}$ is precisely what the streaming interface of
Section~\ref{sec:iostream} avoids unrolling.

\subsection{Long short-term memory}

An LSTM cell~\cite{hochreiter1997} maintains a cell state $c_t$ regulated by
three gates. For input $x_t$ and previous hidden state $h_{t-1}$:

\begin{align}
i_t &= \sigma(W_i x_t + U_i h_{t-1} + b_i) \\
f_t &= \sigma(W_f x_t + U_f h_{t-1} + b_f) \\
o_t &= \sigma(W_o x_t + U_o h_{t-1} + b_o) \\
\tilde{c}_t &= \tanh(W_c x_t + U_c h_{t-1} + b_c) \\
c_t &= f_t \odot c_{t-1} + i_t \odot \tilde{c}_t \\
h_t &= o_t \odot \tanh(c_t)
\end{align}

Two consequences are central to this report. First, the gating requires both a
sigmoid and a hyperbolic tangent, neither of which is cheap in fixed-point
hardware --- both are realised as lookup tables. Second, there are four weight
matrices of each kind, so a cell with $H$ hidden units over $n$ inputs holds
$4H(n + H)$ weights and evaluates $4H$ nonlinearities per time step. The
hardware consequence of that factor of four is measured in
Section~\ref{sec:tables}.

\subsection{Batch normalisation and its folding}

Batch normalisation~\cite{ioffe2015} applies, at inference time, a fixed affine
transform per channel:

\begin{equation}
y = \gamma \cdot \frac{x - \mu}{\sqrt{\sigma^2 + \epsilon}} + \beta .
\end{equation}

Because this is affine and immediately follows a convolution, which is also
affine, the two can be composed into a single convolution with modified weights
and a bias. The synthesis tool performs this automatically, so batch
normalisation costs nothing in hardware. The convolutions in this design are
therefore declared without bias terms, since the folded $\beta$ supplies them.

\subsection{Fixed-point arithmetic}

FPGAs implement floating-point arithmetic only at considerable cost, so
inference is performed in fixed point. The notation
\texttt{ap\_fixed<W,I>} denotes a $W$-bit signed value with $I$ integer bits and
$W - I$ fractional bits, giving a resolution of $2^{-(W-I)}$ and a range of
$\pm 2^{I-1}$. This design uses \texttt{ap\_fixed<16,6>} throughout. The representation, its
rounding and overflow modes, and the automatic widening of intermediate
accumulators are treated in full in Section~\ref{sec:quant}.

\subsection{Reuse factor}

The reuse factor $R$ specifies how many times a single hardware multiplier is
re-used within a layer. For a layer of $N$ MACs, the synthesised design
instantiates $N/R$ multipliers and takes $R$ cycles. Setting $R = N$ yields one
multiplier operating fully sequentially --- minimum area, maximum latency.
Given the 2.56\,s budget, this is the correct operating point for every layer in
this design.

\section{System Description}

\subsection{Overall architecture}

\begin{verbatim}
   MPU6050 (6-axis IMU)                     [not yet integrated]
            |
            v
   Window buffer: 128 samples x 6 channels
            |
            v
   +---------------------------------------+
   |  CNN-LSTM accelerator (FPGA fabric)   |
   |  AXI4-Stream in  : 96 bit  x 128 beats|
   |  AXI4-Stream out : 384 bit x 1 beat   |
   +---------------------------------------+
            |
            v
   argmax -> class -> alert action          [not yet implemented]
\end{verbatim}

\subsection{Network architecture}

\begin{table}[h]
\centering
\begin{tabular}{llr r}
\toprule
\textbf{Layer} & \textbf{Configuration} & \textbf{Output shape} & \textbf{Params} \\
\midrule
Input           & 128 steps $\times$ 6 channels   & (128, 6)  & 0 \\
Conv1D          & 16 filters, $k = 5$, no bias    & (124, 16) & 480 \\
BatchNorm + ReLU& folded into the convolution     & (124, 16) & 64 \\
Conv1D          & 16 filters, $k = 5$, no bias    & (120, 16) & 1{,}280 \\
BatchNorm + ReLU& folded into the convolution     & (120, 16) & 64 \\
MaxPooling1D    & pool 2                          & (60, 16)  & 0 \\
Dropout         & $p = 0.3$, training only        & (60, 16)  & 0 \\
LSTM            & 16 hidden, final state only     & (16,)     & 2{,}176 \\
Dropout         & $p = 0.3$, training only        & (16,)     & 0 \\
Dense + ReLU    & 16 units                        & (16,)     & 272 \\
Dense           & 6 units, logits                 & (6,)      & 102 \\
\midrule
\multicolumn{3}{r}{\textbf{Total trainable parameters}} & \textbf{4{,}374} \\
\bottomrule
\end{tabular}
\caption{Deployed network. Shapes are those reported by the synthesis tool
after batch-normalisation folding and dropout removal.}
\end{table}

The output is six raw logits. The softmax normally applied in training is
omitted in hardware, because it is monotonic and the decision is taken by
\texttt{argmax}, which the normalisation does not affect.

\subsection{Channel selection}

The UCI~HAR dataset provides nine channels: body acceleration, body angular
velocity and total acceleration, each in three axes. This design uses six ---
total acceleration and body angular velocity --- because those are exactly what
an MPU6050 outputs at its register interface. Body acceleration is total
acceleration with the gravity component removed by a Butterworth filter applied
offline by the dataset authors; using it would require reimplementing that
filter in hardware ahead of the network. Discarding it removes that burden
entirely at negligible accuracy cost.

For the same reason no input normalisation is applied. The network is trained
on raw physical units, so on-device preprocessing reduces to a constant scaling
of the sensor's integer output --- a shift rather than a per-channel affine
transform.

\section{Quantization Strategy}
\label{sec:quant}

Neural networks are trained in 32-bit floating point, but floating-point
arithmetic is expensive to synthesise. Deployment therefore requires mapping
every weight, activation and intermediate result onto a fixed-point
representation. This section states what that mapping is in this design, what
the synthesis tool derives automatically, and what the alternatives would be.

\textbf{The quantization used here is post-training and uniform.} The network is
trained entirely in floating point and the trained weights are then mapped to a
fixed-point format specified in the \texttt{hls4ml} configuration. No
quantization-aware training was performed. The consequences of that choice, and
the techniques that would improve on it, are discussed in
Section~\ref{sec:qat}.

\subsection{Uniform fixed-point representation}

A value in the format \texttt{ap\_fixed<W,I>} is a $W$-bit two's-complement
integer interpreted with $I$ bits before the binary point, leaving
$f = W - I$ fractional bits. The representable set is the uniform grid

\begin{equation}
\mathcal{Q}_{W,I} = \left\{ k \cdot \Delta \;\middle|\;
  k \in \mathbb{Z},\;
  -2^{W-1} \le k \le 2^{W-1}-1 \right\},
  \qquad \Delta = 2^{-f},
\end{equation}

so the quantization step and the dynamic range are set independently by $f$ and
$I$ respectively:

\begin{equation}
\Delta = 2^{I-W},
\qquad
x \in \left[\, -2^{\,I-1},\; 2^{\,I-1} - \Delta \,\right].
\end{equation}

This design uses \texttt{ap\_fixed<16,6>} throughout, giving

\begin{equation}
\Delta = 2^{-10} \approx 9.77 \times 10^{-4},
\qquad
x \in [-32,\; 32 - 2^{-10}].
\end{equation}

The format was chosen from the physical units of the input rather than by
search. Accelerometer readings are expressed in $g$ and angular velocity in
rad/s, both of which lie well inside $\pm 32$ for human motion, so six integer
bits provide ample headroom while ten fractional bits resolve the signal far
below sensor noise.

The error introduced by quantization behaves as additive noise of variance
$\Delta^2/12$ for a uniformly distributed residual, giving the standard result
that the signal-to-quantization-noise ratio improves by approximately

\begin{equation}
\Delta \mathrm{SQNR} \approx 20 \log_{10} 2 \approx 6.02\ \mathrm{dB}
\end{equation}

per additional fractional bit. Moving from \texttt{ap\_fixed<16,6>} to
\texttt{ap\_fixed<8,4>} removes six fractional bits and therefore costs roughly
36\,dB of arithmetic fidelity --- which is why the resource saving reported in
Section~\ref{sec:results} cannot be accepted without measuring its accuracy
cost.

\subsection{Rounding and overflow behaviour}

A fixed-point type is incompletely specified by $W$ and $I$ alone: the
behaviour on values falling between grid points, and on values exceeding the
range, must also be fixed. Inspection of the generated project shows the
default modes are in use throughout, the full type being
\texttt{ap\_fixed<16,6,AP\_TRN,AP\_WRAP,0>}.

\textbf{Truncation (\texttt{AP\_TRN}).} The value is rounded toward $-\infty$:

\begin{equation}
Q(x) = \Delta \left\lfloor \frac{x}{\Delta} \right\rfloor,
\qquad
e = Q(x) - x \in (-\Delta,\, 0].
\end{equation}

The error is therefore strictly non-positive, with mean $-\Delta/2$, rather
than the zero-mean error of round-to-nearest (\texttt{AP\_RND}), for which
$e \in (-\Delta/2, \Delta/2]$. This distinction is not cosmetic. Over a dot
product of $N$ terms the truncation bias accumulates \emph{linearly},

\begin{equation}
\mathbb{E}\!\left[\sum_{i=1}^{N} e_i\right] = -\frac{N\Delta}{2},
\end{equation}

whereas a zero-mean rounding error accumulates only as $\mathcal{O}(\sqrt{N})$
in standard deviation. Wider layers are thus systematically biased downward
under truncation. For the 80-term dot products of the second convolutional
layer this is a deterministic offset of $40\Delta \approx 0.039$ before the
nonlinearity.

\textbf{Wrapping (\texttt{AP\_WRAP}).} On overflow the value is reduced modulo
$2^{I}$, so a result marginally above the maximum representable value becomes
the most negative value. Unlike saturation (\texttt{AP\_SAT}), which clamps,
wrapping converts a small numerical excursion into a maximal sign error.
Saturation costs a comparator per value and was not enabled; the design instead
relies on the six integer bits being generous relative to the physical signal
range.

Both defaults are reasonable here given that headroom, but they are defaults
rather than decisions, and they are the first things to revisit should
fixed-point accuracy prove disappointing.

\subsection{Automatic accumulator widening}
\label{sec:accum}

Quantizing weights and activations to $W$ bits does not mean the arithmetic is
performed at $W$ bits. A product of \texttt{ap\_fixed<$W_x$,$I_x$>} and
\texttt{ap\_fixed<$W_w$,$I_w$>} operands requires $W_x + W_w$ bits to be exact,
and summing $N$ such products requires a further $\lceil \log_2 N \rceil$
integer bits to be guaranteed overflow-free. \texttt{hls4ml} sizes these
accumulators automatically, and inspection of the generated type definitions
shows it applies

\begin{align}
W_{\mathrm{acc}} &= W_x + W_w + \lceil \log_2 N \rceil + 1, \\
I_{\mathrm{acc}} &= I_x + I_w + \lceil \log_2 N \rceil + 1,
\end{align}

where the additional bit is a guard bit. The rule was checked against every
layer of the generated design, with $N = k \cdot C_{\mathrm{in}}$ for the
convolutions and $N = n_{\mathrm{in}}$ for the dense layers:

\begin{table}[h]
\centering
\begin{tabular}{lrrll}
\toprule
\textbf{Layer} & \textbf{$N$} & \textbf{$\lceil \log_2 N \rceil$} &
\textbf{Predicted} & \textbf{Generated} \\
\midrule
\texttt{conv1d}    & 30 & 5 & \texttt{ap\_fixed<38,18>} & \texttt{ap\_fixed<38,18>} \\
\texttt{conv1d\_1} & 80 & 7 & \texttt{ap\_fixed<40,20>} & \texttt{ap\_fixed<40,20>} \\
\texttt{dense}     & 16 & 4 & \texttt{ap\_fixed<37,17>} & \texttt{ap\_fixed<37,17>} \\
\texttt{logits}    & 16 & 4 & \texttt{ap\_fixed<37,17>} & \texttt{ap\_fixed<37,17>} \\
\bottomrule
\end{tabular}
\caption{Predicted versus generated accumulator types, for
\texttt{ap\_fixed<16,6>} inputs and weights throughout. The rule reproduces all
four exactly.}
\end{table}

Two practical consequences follow. First, the accumulators are safe by
construction, so overflow cannot occur \emph{inside} a dot product. Second, the
risk is instead concentrated at the layer boundaries, where the wide
accumulator is cast back down to \texttt{ap\_fixed<16,6>} under the wrapping
mode described above. The network output is a visible instance of this: the
logits emerge as \texttt{ap\_fixed<37,17>} rather than the \texttt{<16,6>} used
internally, because the final accumulator is never narrowed. This is harmless
for classification, since \texttt{argmax} is invariant to the representation,
but a host reading those values must interpret them with 17 integer bits.

\subsection{Quantization-aware alternatives}
\label{sec:qat}

Post-training quantization maps a network that was never aware it would be
quantized. Quantization-aware training (QAT) instead simulates the quantizer
during the forward pass, allowing the weights to adapt to it. \textbf{Neither
technique below was used in this project.} Both are the natural next step, and
are described here because they bear directly on the unmeasured accuracy risk
identified in Section~\ref{sec:limitations}.

\textbf{The straight-through estimator.} The obstacle common to all QAT is that
the quantizer is piecewise constant, so $\partial Q / \partial x = 0$ almost
everywhere and no gradient reaches the weights. The standard remedy is the
straight-through estimator, which applies the quantizer in the forward pass but
substitutes the identity in the backward pass:

\begin{equation}
\text{forward: } \hat{x} = Q(x),
\qquad
\text{backward: } \frac{\partial L}{\partial x}
  \;\approx\; \frac{\partial L}{\partial \hat{x}}
  \cdot \mathbf{1}\!\left[\,|x| \le x_{\max}\,\right].
\end{equation}

The indicator suppresses gradients for saturated values, which would otherwise
be pushed further out of range.

\textbf{QKeras} applies this at tensor granularity. Each layer is given a
quantizer such as \texttt{quantized\_bits($b$, $i$, $\alpha$)}, fixing one
bitwidth for an entire weight tensor, and the network is fine-tuned from the
floating-point checkpoint. The bitwidths remain hyperparameters chosen by the
designer.

\textbf{High Granularity Quantization (HGQ)}~\cite{hgq} removes that manual
choice. Rather than one bitwidth per tensor, HGQ assigns a bitwidth $b_i$ to
each individual parameter and learns it jointly with the weights by gradient
descent. For this to be meaningful the hardware cost must enter the objective,
so the training loss is augmented with a differentiable resource proxy. The
natural proxy is the bit-operation count: a multiplication of a $b_w$-bit
weight by a $b_a$-bit activation costs approximately $b_w \cdot b_a$
bit-operations, so for a network with multiplication set $\mathcal{M}$,

\begin{equation}
\mathrm{EBOPs}(\mathbf{b}) \;\approx\;
  \sum_{(w,a) \in \mathcal{M}} b_w \cdot b_a ,
\qquad
\mathcal{L} = \mathcal{L}_{\mathrm{task}}
  + \beta \cdot \mathrm{EBOPs}(\mathbf{b}).
\end{equation}

The bitwidths are held continuous during optimisation and discretised for
export, with the straight-through estimator supplying gradients through the
quantizer. The coefficient $\beta$ traces a Pareto frontier between accuracy
and hardware cost, so instead of selecting a single uniform format the designer
selects an operating point on a measured curve.

This is precisely the choice that was made by hand in this project, and made
crudely: \texttt{ap\_fixed<16,6>} was applied uniformly to every layer, even
though the layers differ substantially in dynamic range, and the 8-bit
alternative was evaluated only as a second uniform format. HGQ would allow the
first convolution, which sees bounded physical units, to use far fewer bits
than the LSTM gates, which do not. Neither library is installed in the current
environment, so adopting either requires toolchain setup in addition to
retraining.

\section{Design and Implementation}

\subsection{Toolflow}

\begin{verbatim}
  Keras / TensorFlow  --->  hls4ml  --->  Vitis HLS  --->  Vivado  ---> FPGA
   train, 90.91%          C++ + tcl        RTL,          IP, block    bitstream
                                           csynth        design
\end{verbatim}

Training uses TensorFlow~2.21 with Adam at a learning rate of $10^{-3}$, batch
size 64, for 40~epochs, with sparse categorical cross-entropy computed from
logits. The dataset split is the one published with UCI~HAR: 7{,}352 training
and 2{,}947 test windows, partitioned by subject so that no participant appears
in both.

\subsection{Framework constraint: PyTorch could not be used}

The network was first implemented in PyTorch. It could not be converted. The
\texttt{hls4ml} PyTorch front end (version~1.3.0) provides no handler for the
\texttt{transpose} operation, and a PyTorch CNN--LSTM requires two of them:
\texttt{nn.Conv1d} expects channel-first tensors while \texttt{nn.LSTM} expects
(batch, time, feature) ordering. Three alternative tensor layouts were
attempted; all failed with \texttt{Unsupported function transpose}.

Keras \texttt{Conv1D} is channels-last by construction, so the identical
architecture converts with no transpose operations at all. The recurrent layer
itself is not the obstacle --- the PyTorch LSTM handler is present and
functional --- which makes this purely an artefact of tensor-layout convention.
The implementation was therefore moved to Keras.

\subsection{Streaming versus parallel interfaces}
\label{sec:iostream}

\texttt{hls4ml} offers two interface models. Under \texttt{io\_parallel} a layer
receives its entire input tensor at once and is unrolled across all output
positions. Under \texttt{io\_stream} the layer receives one time step at a time
through a FIFO and maintains a line buffer, so the hardware instantiated scales
with filter count rather than with window length.

For a convolutional layer producing 124 output positions this distinction is
decisive, and it is not something the reuse factor can compensate for: reuse
shares multipliers \emph{within} a dot product but does not prevent unrolling
\emph{across} output positions. The measured effect is reported in
Section~\ref{sec:results}.

\subsection{Activation layer construction}

Keras offers two syntactically equivalent ways to apply a rectified linear unit:
the \texttt{ReLU} layer class and \texttt{Activation('relu')}. They compute
identical functions, since the \texttt{ReLU} layer's threshold defaults to zero.
They do not synthesise identically. \texttt{hls4ml} maps the layer class to a
\emph{parameterised} thresholded-ReLU primitive that carries a runtime threshold
comparison for every element, while \texttt{Activation('relu')} maps to a plain
ReLU. Across the 3{,}904 activations of the convolutional front end the
difference is substantial, and is quantified in Table~\ref{tab:relu}.

\subsection{Generated IP interface}

Vitis~HLS exports the design as an IP-XACT package for Vivado. The top-level
interface is:

\begin{table}[h]
\centering
\begin{tabular}{lll}
\toprule
\textbf{Port} & \textbf{Width} & \textbf{Description} \\
\midrule
\texttt{imu\_window}  & 96 bit  & AXI4-Stream slave; 6 channels $\times$
                                  \texttt{ap\_fixed<16,6>}, 128 beats \\
\texttt{layer14\_out} & 384 bit & AXI4-Stream master; 6 logits, 1 beat \\
\texttt{ap\_clk}, \texttt{ap\_rst\_n} & 1 bit & clock and active-low reset \\
\texttt{ap\_start}, \texttt{ap\_done} & 1 bit & block-level handshake \\
\texttt{ap\_ready}, \texttt{ap\_idle} & 1 bit & block-level handshake \\
\bottomrule
\end{tabular}
\caption{Top-level ports of the exported IP core.}
\end{table}

Three properties of this interface constrain system integration. Control is a
block-level handshake rather than an AXI4-Lite register map, so a processor
cannot start the core by writing a register; for a streaming kernel
\texttt{ap\_start} is instead tied asserted. The streams carry no
\texttt{TLAST} side-channel, which a DMA engine normally uses to delimit
packets. Finally, 96 and 384 bits are not among the stream widths an AXI~DMA
supports; both are integer multiples of 32~bits, so width converters are
required.

\section{Results}
\label{sec:results}

All figures in this section are Vitis~HLS~2023.2 post-synthesis estimates for
\texttt{xczu7ev-ffvc1156-2-e} unless stated otherwise. Estimates of this kind
are generally pessimistic for LUT count; place-and-route figures are expected to
be lower.

\subsection{Classification accuracy}

\begin{table}[h]
\centering
\begin{tabular}{lrr}
\toprule
\textbf{Configuration} & \textbf{Parameters} & \textbf{Test accuracy} \\
\midrule
16 filters / 16 hidden / 16 dense & 4{,}374  & 90.91\% \\
64 filters / 64 hidden / 64 dense & 60{,}486 & 92.09\% \\
\bottomrule
\end{tabular}
\caption{Accuracy against network width on the UCI~HAR test split.}
\end{table}

A 13.8-fold increase in parameters yields 1.18~percentage points of accuracy.
The smaller network is retained for the reason given in
Section~\ref{sec:width}.

\subsection{Interface model}

\begin{table}[h]
\centering
\begin{tabular}{lrr}
\toprule
\textbf{Block} & \textbf{\texttt{io\_parallel} LUT} & \textbf{\texttt{io\_stream} LUT} \\
\midrule
Conv1D \#1     & 21{,}729  & 2{,}911 \\
Activation \#1 & 113{,}099 & 1{,}075 \\
Conv1D \#2     & 53{,}751  & 4{,}993 \\
Activation \#2 & 109{,}451 & 1{,}107 \\
MaxPooling     & 35        & 1{,}141 \\
\midrule
\textbf{Complete design} & $\approx$\textbf{298{,}000} & \textbf{35{,}745} \\
\bottomrule
\end{tabular}
\caption{Effect of interface model. The \texttt{io\_parallel} column is the
convolutional front end only; the recurrent and dense layers had not been
reported when that synthesis was halted.}
\label{tab:relu}
\end{table}

Two effects are superimposed in this table. The streaming interface removes the
unrolling across output positions. Separately, the activation figures of
113{,}099 and 109{,}451~LUT are an artefact of the thresholded-ReLU construct
described above; replacing it with \texttt{Activation('relu')} reduced those two
blocks to 1{,}075 and 1{,}107~LUT --- a combined reduction from 222{,}550 to
2{,}182~LUT, by a factor of 102, with no change whatsoever to the computed
function or to the 90.91\% accuracy.

\subsection{Target device comparison}

\begin{table}[h]
\centering
\begin{tabular}{lrrrl}
\toprule
\textbf{Device} & \textbf{LUT} & \textbf{DSP} & \textbf{BRAM} & \textbf{Outcome} \\
\midrule
XCZU7EV (ZCU104)   & 16\%  & 3\%   & 25\%  & Fits \\
XC7Z020 (Zynq-7020)& 67\%  & 24\%  & 56\%  & Would fit \\
GW1NR-9 (Tang Nano 9K) & 414\% & 270\% & 612\% & Does not fit \\
\bottomrule
\end{tabular}
\caption{Utilisation of the same 35{,}745~LUT design across three candidate
devices.}
\end{table}

\subsection{Final design on the target device}

\begin{table}[h]
\centering
\begin{tabular}{lrr}
\toprule
\textbf{Resource} & \textbf{Used} & \textbf{Utilisation} \\
\midrule
LUT              & 33{,}500 & 14\% \\
Flip-flops       & 21{,}724 & 4\%  \\
DSP slices       & 54       & 3\%  \\
Block RAM (18\,Kb) & 159    & 25\% \\
UltraRAM         & 0        & 0\%  \\
\midrule
Clock            & \multicolumn{2}{l}{100\,MHz, timing met (0.37\,ns slack)} \\
Latency          & \multicolumn{2}{l}{160{,}339 cycles = 1.60\,ms} \\
Throughput margin& \multicolumn{2}{l}{2.56\,s window / 1.60\,ms = 1600$\times$} \\
\bottomrule
\end{tabular}
\caption{Final configuration: \texttt{io\_stream}, maximum reuse,
\texttt{ap\_fixed<16,6>}, on XCZU7EV at 100\,MHz.}
\end{table}

\subsection{Arithmetic precision}

\begin{table}[h]
\centering
\begin{tabular}{lrrr}
\toprule
\textbf{Precision} & \textbf{LUT} & \textbf{DSP} & \textbf{BRAM} \\
\midrule
\texttt{ap\_fixed<16,6>} & 35{,}745 & 54 & 159 \\
\texttt{ap\_fixed<8,4>}  & 26{,}526 & 16 & 109 \\
\bottomrule
\end{tabular}
\caption{Effect of halving the word length (50\,MHz runs).}
\end{table}

Halving the word length reduces LUT usage by 26\%, block RAM by 31\% and DSP
usage by 70\%. The accuracy cost of this reduction has not been measured, and
16-bit precision was retained because the target device has ample margin.

\subsection{Activation lookup tables}
\label{sec:tables}

The LSTM's sigmoid and tanh evaluations are realised as lookup tables, and
these dominate the recurrent layer's memory footprint. Two experiments isolate
what drives their size.

\begin{table}[h]
\centering
\begin{tabular}{lrr}
\toprule
\textbf{Table} & \textbf{$H = 16$} & \textbf{$H = 64$} \\
\midrule
sigmoid & 24 BRAM & 96 BRAM \\
tanh    & 8 BRAM  & 32 BRAM \\
\midrule
Total   & 32 BRAM & 128 BRAM \\
\bottomrule
\end{tabular}
\caption{Activation-table memory against LSTM hidden size.}
\end{table}

Reducing the configured table depth from 1024 to 256 entries changed nothing:
the generated constant was verified to be 256 in the synthesised project, and
the reported resources were byte-identical. Table footprint therefore does not
scale with depth. It scales with hidden size, exactly linearly --- a factor of
four in $H$ producing a factor of four in memory --- which is consistent with
the table being replicated to permit parallel access across the $4H$ gate
evaluations each time step requires.

\subsection{Network width and the memory ceiling}
\label{sec:width}

\begin{table}[h]
\centering
\begin{tabular}{lrrrr}
\toprule
\textbf{Configuration} & \textbf{LUT} & \textbf{DSP} & \textbf{BRAM} & \textbf{Latency} \\
\midrule
16/16/16 & 33{,}500 (14\%)  & 54 (3\%)   & 159 (25\%)            & 1.60\,ms \\
64/64/64 & 126{,}167 (54\%) & 198 (11\%) & \textbf{654 (104\%)}  & 25.4\,ms \\
\bottomrule
\end{tabular}
\caption{The wider network exceeds the device's block~RAM.}
\end{table}

The wider network is comfortable in logic and arithmetic but exceeds available
block~RAM by 4\%. Since the device's 96~UltraRAM blocks (27\,Mb) remain entirely
unused, this is a mapping problem rather than a capacity one, and directing the
weight arrays to UltraRAM is the obvious remedy. That work was not completed, so
the narrower network remains the only configuration measured to fit, and the
1.18-point accuracy difference stands as its cost.

\subsection{Negative result: infeasibility on a small FPGA}

The project initially targeted a Sipeed Tang~Nano~9K (Gowin GW1NR-9: 8{,}640
LUT4, 20~DSP, 26 block~RAMs). The design does not fit, and the reason is
structural rather than a question of tuning.

The LSTM accounts for 56\% of all logic in the design. Within it, the sigmoid
and tanh tables alone require 32~block~RAMs --- more than the 26 the entire
device contains --- before a single weight or stream buffer is allocated. Since
gated recurrence requires both nonlinearities by definition, no reduction in
precision, table depth or weight count removes this floor. Measurements
confirmed this: 8-bit precision brought DSP usage within budget (16 of 20) but
left LUT usage 3.1$\times$ and block~RAM 4.2$\times$ over.

The practical conclusion is that an LSTM is the wrong primitive for a
sub-10k-LUT device. A purely convolutional network with global pooling has no
such tables; the measured convolutional front end occupies 7{,}904~LUT and would
fit that device with room remaining for peripheral logic.

\section{Applications and Advantages}

\begin{itemize}
  \item Real-time activity classification with a 1600$\times$ latency margin,
        leaving substantial headroom for additional processing stages.
  \item Inference performed locally, so raw inertial data need never leave the
        device --- relevant where the sensor stream is privacy-sensitive.
  \item A reproducible toolflow in which retraining the network and
        regenerating hardware requires no manual RTL modification.
  \item Deterministic latency, as the design contains no cache, operating
        system or contended shared resource.
  \item A characterised design space, allowing the accuracy/area trade-off to be
        re-evaluated quickly if the target device changes.
\end{itemize}

\section{Limitations and Future Scope}
\label{sec:limitations}

\subsection{Limitations}

\begin{itemize}
  \item \textbf{No hardware deployment.} All results are high-level-synthesis
        estimates. No bitstream has been generated and the design has not been
        executed on a physical device. Place-and-route figures will differ.
  \item \textbf{Fixed-point accuracy unverified.} The 90.91\% figure is the
        floating-point model's. The accuracy of the 16-bit fixed-point
        implementation has not been measured. Since the quantization is
        post-training and uniform, with no quantization-aware fine-tuning,
        there is no guarantee the mapping is benign.
  \item \textbf{The sensor is not integrated.} The MPU6050 interface, window
        assembly and alert logic are specified but unimplemented.
  \item \textbf{The task is activity recognition, not safety detection.}
        UCI~HAR contains six ordinary activities and no fall, impact or
        distress class. The trained network cannot detect events it has never
        observed, and no claim of fall or assault detection is supported by
        these results.
  \item \textbf{Window-level, not event-level.} Classification is performed on
        2.56\,s windows in isolation, with no temporal smoothing or
        hysteresis across consecutive windows.
\end{itemize}

\subsection{Future scope}

\begin{itemize}
  \item Generate the bitstream and validate on hardware, comparing
        place-and-route utilisation against the estimates reported here.
  \item Measure fixed-point accuracy using bit-accurate C simulation, and
        re-evaluate 8-bit precision against the measured loss.
  \item Replace uniform post-training quantization with quantization-aware
        training. High Granularity Quantization (Section~\ref{sec:qat}) is the
        strongest candidate, since it learns a per-parameter bitwidth against a
        differentiable resource cost rather than requiring one format to be
        chosen by hand for the whole network.
  \item Map weight arrays to UltraRAM to accommodate the wider network and
        recover the 1.18-point accuracy difference.
  \item Retrain on a fall-annotated dataset (for instance EdgeFall) to support
        the safety application the architecture was originally intended for.
  \item Integrate the MPU6050 and implement continuous windowed inference with
        temporal smoothing.
  \item Measure power consumption, which is the metric that would justify an
        FPGA over a microcontroller for a battery-powered device, and which
        this work does not address.
\end{itemize}

\section{Conclusion}

A 1D~CNN--LSTM activity classifier was trained to 90.91\% accuracy on UCI~HAR
using 4{,}374 parameters, converted to synthesisable hardware through
\texttt{hls4ml} and Vitis~HLS, and packaged as an AXI4-Stream IP core. On the
XCZU7EV the design occupies 14\% of available logic and completes inference in
1.60\,ms against a 2.56\,s window.

The more transferable finding concerns how little of the hardware cost was
determined by the network itself. Two configuration choices --- the activation
layer construct and the interface model --- together accounted for roughly an
order of magnitude in area, while leaving the computed function unchanged. The
corresponding negative result is equally concrete: the sigmoid and tanh tables
an LSTM requires exceed the entire block~RAM capacity of a small FPGA
independently of precision or table depth, which rules out gated recurrence on
such devices rather than merely making it difficult.

The flow from trained network to packaged IP is complete and reproducible. What
remains is the hardware half: bitstream generation, on-board validation, sensor
integration, and retraining on data containing the events a safety device would
need to recognise.

\section{References}

\begin{thebibliography}{99}

\bibitem{anguita2013}
Anguita, D., Ghio, A., Oneto, L., Parra, X., \& Reyes-Ortiz, J.~L. (2013).
\textit{A Public Domain Dataset for Human Activity Recognition Using
Smartphones}. 21st European Symposium on Artificial Neural Networks (ESANN).

\bibitem{ordonez2016}
Ordóñez, F.~J., \& Roggen, D. (2016). \textit{Deep Convolutional and LSTM
Recurrent Neural Networks for Multimodal Wearable Activity Recognition}.
Sensors, 16(1), 115.

\bibitem{hammerla2016}
Hammerla, N.~Y., Halloran, S., \& Plötz, T. (2016). \textit{Deep,
Convolutional, and Recurrent Models for Human Activity Recognition using
Wearables}. Proceedings of IJCAI 2016.

\bibitem{duarte2018}
Duarte, J., et al. (2018). \textit{Fast inference of deep neural networks in
FPGAs for particle physics}. Journal of Instrumentation, 13, P07027.

\bibitem{umuroglu2017}
Umuroglu, Y., Fraser, N.~J., Gambardella, G., Blott, M., Leong, P., Jahre, M.,
\& Vissers, K. (2017). \textit{FINN: A Framework for Fast, Scalable Binarized
Neural Network Inference}. Proceedings of FPGA 2017.

\bibitem{hochreiter1997}
Hochreiter, S., \& Schmidhuber, J. (1997). \textit{Long Short-Term Memory}.
Neural Computation, 9(8), 1735--1780.

\bibitem{ioffe2015}
Ioffe, S., \& Szegedy, C. (2015). \textit{Batch Normalization: Accelerating
Deep Network Training by Reducing Internal Covariate Shift}. Proceedings of
ICML 2015.

\bibitem{lecun2022}
LeCun, Y. (2022). \textit{A Path Towards Autonomous Machine Intelligence}.
OpenReview.

\bibitem{assran2023}
Assran, M., et al. (2023). \textit{Self-Supervised Learning from Images with a
Joint-Embedding Predictive Architecture}. Proceedings of CVPR 2023.

\bibitem{hgq}
Sun, C., \AA{}rrestad, T.~K., Loncar, V., Ngadiuba, J., \& Spiropulu, M.
(2024). \textit{Gradient-based Automatic Mixed Precision Quantization for
Neural Networks On-Chip}. (High Granularity Quantization, HGQ.)
% TODO verify: venue, year and author list before submission.

\bibitem{vitishls}
AMD/Xilinx. \textit{Vitis High-Level Synthesis User Guide} (UG1399),
version 2023.2.

\bibitem{zcu104}
AMD/Xilinx. \textit{ZCU104 Evaluation Board User Guide} (UG1267).

\end{thebibliography}

\newpage
\appendix

\section{Development Environment}

\begin{table}[h]
\centering
\begin{tabular}{lll}
\toprule
\textbf{Component} & \textbf{Version} & \textbf{Role} \\
\midrule
Python            & 3.11.9      & host environment \\
TensorFlow / Keras& 2.21 / 3.14 & network training \\
\texttt{hls4ml}   & 1.3.0       & network $\rightarrow$ HLS C++ \\
Vitis HLS         & 2023.2      & C++ $\rightarrow$ RTL, IP packaging \\
Vivado            & 2023.2      & block design, implementation \\
Target device     & \multicolumn{2}{l}{XCZU7EV-2FFVC1156E (ZCU104)} \\
\bottomrule
\end{tabular}
\end{table}

\section{Target Device Resources}

\begin{table}[h]
\centering
\begin{tabular}{lrr}
\toprule
\textbf{Resource} & \textbf{XCZU7EV} & \textbf{GW1NR-9} \\
\midrule
LUT             & 230{,}400 & 8{,}640 \\
Flip-flops      & 460{,}800 & --- \\
DSP slices      & 1{,}728   & 20 \\
Block RAM       & 11.0\,Mb  & 468\,Kb \\
UltraRAM        & 27.0\,Mb  & --- \\
Processor       & \multicolumn{1}{c}{4$\times$Cortex-A53} & none \\
\bottomrule
\end{tabular}
\end{table}

\section{Reproduction}

\begin{lstlisting}[language=bash]
# Build the dataset cache and train the network
python -m jepasafe.demopart --epochs 30
python -m jepasafe.cnnlstm_keras --epochs 40

# Convert to an HLS project
python -m jepasafe.export_hls

# C synthesis
export PATH="/c/Xilinx/Vitis_HLS/2023.2/bin:$PATH"
cd hls/jepasafe && vitis_hls -f build_prj.tcl
# report: jepasafe_prj/solution1/syn/report/csynth.rpt

# Package as a Vivado IP
cd hls/jepasafe && vitis_hls -f ../../scripts/export_ip.tcl
# output: jepasafe_prj/solution1/impl/ip/nitt_hls_jepasafe_1_0.zip
\end{lstlisting}

\section{Dataset Summary}

\begin{table}[h]
\centering
\begin{tabular}{ll}
\toprule
Source           & UCI HAR (Anguita et al., 2013) \\
Participants     & 30, split by subject between train and test \\
Sampling rate    & 50\,Hz \\
Window           & 128 samples (2.56\,s) \\
Training windows & 7{,}352 \\
Test windows     & 2{,}947 \\
Channels used    & 6 of 9 (total acceleration, body angular velocity) \\
Classes          & walking, upstairs, downstairs, sitting, standing, laying \\
\bottomrule
\end{tabular}
\end{table}

\end{document}
